% DEFINED:   sec:rnn-pattern, sec:rnn-handoff
% RESOLVED:  ch:cnn, sec:mlp-backprop (backward into Part II and Part I)
% FORWARD (prose-only): Chapter 6 (the Transformer, "the next chapter").
%            Not drafted yet; stays prose, no \Cref.

\section{The pattern repeats}\label{sec:rnn-pattern}

Three architectures now sit on the table, and they line up along one idea.

\begin{center}
\begin{tabular}{@{}lll@{}}
\toprule
Architecture & Weight-sharing axis & Equivariance \\
\midrule
MLP & none & input permutation \\
CNN & spatial position & 2-D translation in space \\
RNN & timestep & 1-D translation in time \\
\bottomrule
\end{tabular}
\end{center}

The recipe is the same each time. Identify a symmetry the data respects.
Wire it into the architecture by reusing one set of weights along the axis
the symmetry runs along. Three things follow at once. The parameter count
drops, because one shared operator replaces a separate operator per
position. The data efficiency rises, because what the model learns at one
position transfers to all the others for free. And the backward pass becomes
a \texttt{+=} accumulation along the shared axis, the pattern the library has
carried since the \texttt{Linear} layer of \Cref{sec:mlp-backprop}: across
examples in a mini-batch, then across spatial positions in a convolution of
\Cref{ch:cnn}, then across timesteps in a recurrent unroll. One operation,
three grains.

What the recurrent network buys over a fixed-context model, for sequence
data, is the state. Information can in principle flow from an arbitrarily
distant past input through the hidden vector, where a fixed window simply
cannot reach past its edge. What it costs is the vanishing-gradient problem.
The state is a single fixed-width vector that every past input has to pass
through, and the gradient that would teach the model to use a distant input
has to survive a long product of Jacobians to get back there. The reach is
unbounded in principle and limited in practice, and the limit is structural,
not a matter of more training.

\section{Handoff to the Transformer}\label{sec:rnn-handoff}

The next chapter takes the other path. Rather than squeezing the past
through a sequential state, the Transformer lets every position attend to
every other position directly, in a single layer. That one move changes
three things at once. It removes the long Jacobian product, because there is
no recurrence to backpropagate through, so the vanishing-gradient problem of
this chapter goes away. It restores parallelism over the sequence, because
without a state to thread, the positions no longer have to be processed in
order. And it trades time-translation equivariance for permutation
equivariance over positions, which then has to be partly broken back with a
positional encoding so the model can still tell ``the cat ate the mouse''
from ``the mouse ate the cat.''

That is a different bet about which symmetries a sequence respects. The
recurrent network commits to time-translation equivariance and a running
state, and pays with sequential training and a fragile gradient. The
Transformer commits to permutation equivariance, breaks it deliberately with
a positional encoding, and pays with a cost that grows quadratically in the
sequence length. Both are entries on the same menu of architectural priors,
and the right one is the one whose assumptions match the data. The next
chapter builds the Transformer from the same pure-derivation, real-code
footing this one used, and asks what content-addressable attention buys that
a state cannot.
